% !PRINTED 406
\weberpage{406}
%<<CONT
worin $m$ alle positiven ganzen Zahlen, die zu $Q$ relativ prim sind,
und $n$ für jedes $m$ die Teiler von $m$ durchläuft.

Ordnen wir auf der rechten Seite von (2) nach $n$, so kommen
für ein gegebenes $n$ unter den $m$ alle Vielfachen von $n$ vor, und
wir können auch setzen:

\weq{3}{\Sigma F(N\mathfrak{a}) = \overset{m}{\Sigma}\overset{n}{\Sigma}(\Delta,n)\,F(mn),}

worin jetzt $m$ und $n$, voneinander unabhängig, alle natürlichen
Zahlenwerte annehmen, die mit $Q$ keinen Teiler gemein haben.

Ist im besonderen $F$ so beschaffen, daß

\weq{4}{F(mn) = F(m)\,F(n)}

ist, so können wir (3) auch so darstellen:

\weq{5}{\overset{\mathfrak{a}}{\Sigma}F(N\mathfrak{a}) = \overset{m}{\Sigma}F(m)\overset{n}{\Sigma}(\Delta,n)\,F(n).}

Bei den Summen nach $m$ und $n$ sind die Zahlen ausgeschlossen
, die mit $Q$ einen Teiler gemein haben.

Ebenso sind auf der linken Seite die Ideale ausgeschlossen,
die zu $Q$ nicht relativ prim sind.

Die Formel (5) können wir zunächst anwenden, um aus den
Formeln des vorigen Paragraphen die \gesperrt{Klassenzahl} zu bestimmen.
Man erhält auf diese Weise auch das Verhältnis der Klassenzahl
im \gesperrt{Körper} zu der Klassenzahl nach der Ordnung $[Q]$. Da wir
dieses Verhältnis aber schon auf anderem Wege bestimmt haben
(§ 100), so wollen wir uns hier auf die \gesperrt{Klassenzahl des Körpers}
 beschränken, d.\,h. wir wollen $Q = 1$ setzen.

Setzen wir in (5) $F(z) = z^{-s}$, so findet die vorausgesetzte
unbedingte Konvergenz statt, solange $s > 1$ ist. Multiplizieren
wir aber mit $s - 1$ und lassen $s$ in $1$ übergehen, so können wir
von § 111, (12) Gebrauch machen.

Diese Formel wenden wir in der ersten Fassung an, bei
der $\mathfrak{m}$ die zu $Q$ teilerfremden Ideale der Klasse $M$ durchläuft
und $r$ die Primfaktoren von $Q$ bedeuten. Diese fallen also für
$Q = 1$ ganz weg, und $\mathfrak{m}$ ist in der Summenformel:

\weq{6}{\begin{aligned}
\sum\frac{s-1}{(N\mathfrak{m})^s} &= \frac{1}{\sqrt{\Delta}}\log\varepsilon, & \Delta&>0\\
&= \frac{\pi}{\sqrt{-\Delta}}, & \Delta&<0
\end{aligned}}

keiner Beschränkung in bezug auf Teilbarkeit mehr unterworfen
.
